\documentclass[11pt]{article}
\usepackage{amsmath,amssymb,amsthm,array}
\usepackage[margin=1in]{geometry}
\usepackage[hidelinks]{hyperref}

\newcommand{\CP}{\mathbb{CP}}
\newcommand{\RR}{\mathbb{R}}
\newcommand{\OO}{\mathcal{O}}
\newcommand{\WW}{\mathcal{W}}
\newcommand{\ii}{\mathrm{i}}
\newcommand{\Imop}{\operatorname{Im}}
\newcommand{\Vol}{\operatorname{Vol}}
\newcommand{\eff}{\mathrm{eff}}

\theoremstyle{plain}
\newtheorem{proposition}{Proposition}[section]
\newtheorem{lemma}[proposition]{Lemma}
\theoremstyle{definition}
\newtheorem{example}[proposition]{Example}

\title{D10 boundary-current advances for the de Vries slice}
\author{dv2 working note}
\date{2026-05-25}

\begin{document}
\maketitle

\begin{abstract}
This note collects the new calculations in the D10 boundary-current route to
the de Vries electroweak number.  The construction uses Type IIA data on
\(M_4\times\CP^2\times\CP^1\), a D0/M-circle probe side, and a D8-like
boundary carrier \(\WW_9=M_4\times\CP^2\times S^1_Q\).  It gives a
positive image-channel Feshbach algebra, gives \(\kappa_{\rm cur}=1\) on the
full-current branch, and packages the scalar slice in the
coupled variables \(\rho,\tau,v_Q,\lambda_H,\kappa_{\rm cur}\).  The new
residue calculation shows that the raw Wilson second variation has residue
\(J_j\), while the finite Feshbach pole requires unit residue.  The unit inverse
image metric is the static Schur complement of the retained image oscillator.
For a generic local oscillator the pole depends on a residue ratio
\(r=\gamma/\alpha\).  A single current-displacement square forces \(r=1\).
The mixed Chan--Paton sector \(H^0(\CP^1,\operatorname{Hom}(\OO(0),\OO(2j)))\)
has the right rank: \(Q_j=A_jy_j/\sqrt{J_j}\) lies in \(\Im A_j\) for
\(j>0\).  The current D10 boundary package contains no one-square Hom action
that would lock its coefficient.
The D8/Payen/RR source package also contains no term fixing the full-current placement:
\(\kappa_{\rm cur}=1\) is a branch condition in the written action, while a
horizontal-only image channel gives \(\sin^2\theta_W=0.242265\ldots\).
The neutral scalar Hessian has enough coefficient freedom to fit the
negative-branch scalar constants; it does not force them in the present action.
The source-balance equations reduce any replacement source sector to a
two-component lower-charge target; the minimal target is \(-h_Q^2\), not a
field already present in the current D8/O8 package.
A tower/moduli branch instead retains the \(j=3/2\) state, with mass
\(96.54\,{\rm GeV}\) when the scale is fixed by \(M_Z\).
\end{abstract}

% The local string derivation of the inverse-image-metric contact is not printed
% as a result here, because it has not been derived from the boundary action.

\section{The de Vries datum}

Let
\begin{equation}
  X_+(J)=\frac{-J+\sqrt{J^2+4J}}{2},
  \qquad J=j(j+1).
  \label{eq:xplus}
\end{equation}
The two electroweak slots use
\begin{equation}
  J_W=\frac34,\qquad J_Z=2.
\end{equation}
Then
\begin{equation}
  \sin^2\theta_{\rm dV}
  =1-\frac{X_+(3/4)}{X_+(2)}
  =0.2231013223008662,
  \label{eq:dv-sin}
\end{equation}
and
\begin{equation}
  \frac{M_W}{M_Z}
  =\sqrt{\frac{X_+(3/4)}{X_+(2)}}
  =0.8814185598789793.
  \label{eq:dv-ratio}
\end{equation}
The D10 calculation below asks whether this number is forced by local boundary
dynamics, not only placed in a compatible parameter slice.

\section{D10 carrier and probe story}

The active ten-dimensional layer is Type IIA on
\begin{equation}
  M_4\times K_6,\qquad K_6=\CP^2\times\CP^1 .
  \label{eq:d10}
\end{equation}
The M-circle reduction gives the string-frame form
\begin{equation}
  ds_{11}^2=e^{-2\Phi/3}ds_{10,\rm str}^2
  +e^{4\Phi/3}(d\psi+C_1)^2,
  \qquad F_2=dC_1 .
  \label{eq:m-circle}
\end{equation}
The D0-brane is the electric probe of \(C_1\):
\begin{equation}
  S_{\rm D0,WZ}=\mu_0\int_{W_1}C_1,
  \qquad Q_{\rm D0}\leftrightarrow P_\psi .
  \label{eq:d0}
\end{equation}
The space-filling side in Type IIA is the D8-like boundary carrier
\begin{equation}
  \WW_9=M_4\times\CP^2\times S^1_Q.
  \label{eq:w9}
\end{equation}
It is codimension one in \(K_6\).  The literal all-Neumann D9 row is retained
only as an open-string control row.

\begin{proposition}[Carrier/probe separation]
The D0 and \(\WW_9\) rows test different pieces of one D10 background.  The
D0 measures the M-circle/RR one-form.  The D8-like boundary carries the
\(Q\)-preserving endpoint/current sector and can support the Chan--Paton
representations used by the weak boundary modes.
\end{proposition}

\begin{proof}
Equation \eqref{eq:d0} couples the D0 worldline to \(C_1\), the connection in
\eqref{eq:m-circle}.  Equation \eqref{eq:w9} wraps the observed spacetime,
the color \(\CP^2\), and the weak \(Q\)-circle, leaving one weak normal
direction.  The worldvolume dimension and RR parity
therefore match a Type IIA D8-like source, not a D9-brane.
\end{proof}

\section{Boundary action and photon protection}

The weak endpoint spaces are
\begin{equation}
  V_j=H^0(\CP^1,\OO(2j)),\qquad j\in \tfrac12\mathbb Z_{\ge0}.
  \label{eq:vj}
\end{equation}
Borel--Weil gives \(T^aT_a=j(j+1)\mathbf 1_{V_j}\).  The finite current branch
keeps the level-two set
\begin{equation}
  j=0,\ \frac12,\ 1.
  \label{eq:level-two}
\end{equation}
The tower branch keeps higher \(j\).

The \(Q\)-preserving boundary field is a section
\begin{equation}
  \eta\in H^0(\CP^1,\OO(1))\otimes L_Y,
  \qquad Q=T_3+Y.
  \label{eq:eta}
\end{equation}
Choosing the lower doublet component \(\eta_0=v\) gives
\begin{equation}
  T_3\eta_0=-\frac12\eta_0,\qquad
  Y\eta_0=\frac12\eta_0,\qquad
  Q\eta_0=0.
  \label{eq:qeta}
\end{equation}

\begin{proposition}[Photon protection]
On the \(\eta=\eta_0\) boundary slice, the \(Q\)-connection has no mass term
from \(\langle D\eta,D\eta\rangle\).
\end{proposition}

\begin{proof}
The contribution proportional to the \(Q\)-connection is
\(\ii A_QQ\eta_0\), which vanishes by \eqref{eq:qeta}.  Generators that do not
annihilate \(\eta_0\) contribute quadratic terms
\(\langle R\eta_0,R\eta_0\rangle A_R^2\).
\end{proof}

\section{Image-channel Feshbach algebra}

Let \(e_a\), \(a=1,2,3\), be an orthonormal weak-current frame.  Define
\begin{equation}
  A_j:V_j\longrightarrow V_j\otimes\RR^3,\qquad
  A_jp=\sum_{a=1}^3T_a^{(j)}p\otimes e_a .
  \label{eq:aj}
\end{equation}
Then
\begin{equation}
  A_j^\dagger A_j=J_j\mathbf 1_{V_j},
  \qquad J_j=j(j+1).
  \label{eq:ajnorm}
\end{equation}
For \(j>0\), let \(q_j\in\Imop A_j\) and
\begin{equation}
  P_j=\frac{1}{J_j}A_jA_j^\dagger
  \label{eq:proj}
\end{equation}
be the image projector.

The positive image-channel quadratic form is
\begin{equation}
\begin{split}
  S_{\rm quad}^{(j)}
  =
  \langle p_j,Bp_j\rangle
  +\langle q_j,(B+\mu^2J_j)q_j\rangle
  -\mu^2\langle q_j,A_jp_j\rangle
  -\mu^2\langle A_jp_j,q_j\rangle .
  \label{eq:squad}
\end{split}
\end{equation}
Eliminating \(q_j\) gives
\begin{equation}
  \Gamma_j(B)
  =
  B-\mu^4A_j^\dagger(B+\mu^2J_j)^{-1}A_j
  =
  B-\frac{\mu^4J_j}{B+\mu^2J_j}.
  \label{eq:gamma}
\end{equation}
With \(X=B/\mu^2\), the pole equation is
\begin{equation}
  X=\frac{J_j}{X+J_j},
  \label{eq:feshbach-fixed}
\end{equation}
which is \eqref{eq:xplus}.

\section{Canonical current metric}

Let \(K_a\) be the weak Killing/current generators in the same internal metric
that appears in the Kaluza--Klein reduction.  Define
\begin{equation}
  G^{(2)}_{ab}
  =
  \frac{1}{\Vol(\widehat K_6)}
  \int_{\widehat K_6}\sqrt g\,g_{mn}K_a^mK_b^n .
  \label{eq:metric}
\end{equation}
On the \(SU(2)\)-symmetric weak line,
\begin{equation}
  G^{(2)}_{ab}=L_2^2\delta_{ab},
  \qquad e_a=\frac{K_a}{L_2}.
  \label{eq:frame}
\end{equation}
In this frame \eqref{eq:ajnorm} gives
\begin{equation}
  \kappa_{\rm cur}=1.
  \label{eq:kappaone}
\end{equation}

This conclusion is conditional on the boundary image channel using the full
current \(K_{a,\rm full}\).  The Bailin--Love split has
\begin{equation}
  L_{2,BL}^2=L_{2,horiz}^2+L_{2,vert}^2
  =(1+\beta)L_{2,horiz}^2 .
  \label{eq:blsplit}
\end{equation}
If the endpoint image channel sees only the horizontal \(\CP^1\) current, then
\begin{equation}
  \kappa_h=\frac{L_{2,horiz}^2}{L_{2,BL}^2}
  =\frac{1}{1+\beta}=0.7527390090748721,
  \label{eq:kappah}
\end{equation}
and the deformed pole gives
\begin{equation}
  \sin^2\theta_W=0.24226537471822063.
\end{equation}
The Payen endpoint action supplies the representation tensor for the
connection inserted into the boundary Wilson loop.  The D8 WZ term supplies RR
charge data.  Neither term in the current source package fixes the placement
of \(L_{2,vert}\) inside the endpoint current metric.

\begin{proposition}[Full-current normalization]
The boundary image map and the weak gauge metric share the same Casimir
normalization on the full-current branch.  A horizontal-only boundary current
changes the normalization to
\[
  \kappa_{\rm cur}=0.7527390090748721,
  \qquad
  \sin^2\theta=0.24226537471822063.
\]
\end{proposition}

\begin{proof}
Equation \eqref{eq:frame} orthonormalizes the same current vectors used in the
Kaluza--Klein gauge metric.  Hence
\[
  A_j^\dagger A_j
  =\sum_{a,b}T_a^{(j)}T_b^{(j)}\langle e_a,e_b\rangle
  =\sum_aT_a^{(j)}T_a^{(j)}
  =J_j\mathbf 1_{V_j}.
\]
Replacing the boundary image metric by only the horizontal component rescales
the current frame by the recorded factor, which shifts the weak-angle output.
\end{proof}

\section{Wilson-contact residue}

The static part of \eqref{eq:squad} is
\begin{equation}
  S_{\rm stat}
  =
  \mu^2J_j\|q_j\|^2
  -\mu^2\langle q_j,A_jp_j\rangle
  -\mu^2\langle A_jp_j,q_j\rangle .
  \label{eq:stat}
\end{equation}
Completing the square in \(\Imop A_j\) gives
\begin{equation}
  S_{\rm stat}
  =
  \mu^2J_j\left\|q_j-\frac{A_jp_j}{J_j}\right\|^2
  -\mu^2\|p_j\|^2.
  \label{eq:contact-fixed}
\end{equation}
Thus the finite image-channel pole fixes the open-channel contact as
\begin{equation}
  C_{\rm req}=-\mu^2\mathbf 1_{V_j}.
  \label{eq:creq}
\end{equation}

The Wilson factor expansion gives
\begin{equation}
  \frac{\delta^2W_j}{\delta\xi^a\delta\xi^b}\bigg|_{\xi=0}
  \sim -\mu^2T_a^{(j)}T_b^{(j)}.
  \label{eq:wilson2}
\end{equation}
In the orthonormal current frame this contracts to
\begin{equation}
  C_{\rm Wilson}
  =
  -\mu^2\sum_aT_a^{(j)}T_a^{(j)}
  =
  -\mu^2J_j\mathbf 1_{V_j}.
  \label{eq:cwilson}
\end{equation}
Projection to the image does not change the coefficient:
\begin{equation}
  A_j^\dagger P_jA_j
  =
  A_j^\dagger A_j
  =
  J_j\mathbf 1_{V_j}.
  \label{eq:projection-nochange}
\end{equation}

\begin{proposition}[Residue test]
The raw Wilson second variation has residue \(J_j\), while the finite
Feshbach pole requires unit residue.  The unit contact is
\begin{equation}
  -\mu^2A_j^\dagger
  \left(A_jA_j^\dagger|_{\Imop A_j}\right)^{-1}
  A_j
  =
  -\mu^2\mathbf 1_{V_j}.
  \label{eq:inverse-image}
\end{equation}
\end{proposition}

\begin{proof}
Equations \eqref{eq:contact-fixed} and \eqref{eq:creq} fix the contact required
by the pole.  Equations \eqref{eq:cwilson} and \eqref{eq:projection-nochange}
compute the contact from the Wilson Taylor term and from the projected image
tensor.  On \(\Imop A_j\),
\(A_jA_j^\dagger=J_j\mathbf 1\), so
\[
  A_j^\dagger
  \left(A_jA_j^\dagger|_{\Imop A_j}\right)^{-1}
  A_j
  =
  \frac{1}{J_j}A_j^\dagger A_j
  =
  \mathbf 1_{V_j}.
\]
\end{proof}

If the inverse image metric is absent and the Wilson residue is used, the
kernel becomes
\begin{equation}
  \Gamma_j(B)
  =
  B+\mu^2(1-J_j)
  -\frac{\mu^4J_j}{B+\mu^2J_j}.
  \label{eq:shifted}
\end{equation}
The pole equation is
\begin{equation}
  X+1-J_j-\frac{J_j}{X+J_j}=0.
  \label{eq:shifted-pole}
\end{equation}
This changes the W and Z slots by different amounts and cannot be absorbed
into one common scale.

The inverse in \eqref{eq:inverse-image} is not a separate microscopic operator
when \(q_j\) is retained.  It is the \(B=0\) Schur complement of the local
image oscillator:
\begin{equation}
  \Gamma_j(0)
  =
  -\mu^2A_j^\dagger
  \left(A_jA_j^\dagger|_{\Imop A_j}\right)^{-1}
  A_j
  =
  -\mu^2\mathbf 1_{V_j}.
  \label{eq:static-schur}
\end{equation}
The finite branch is shifted only if the raw Wilson \(T_aT_b\) term is also
kept as a direct 1PI \(p_j\) contact.

The endpoint action is linear in the source.  Therefore the Wilson second
derivative is a connected two-current insertion rather than a microscopic
quadratic seagull.  If \(q_j\) is retained as the local image-current channel,
the connected response is represented by \(q_j\) exchange and should not be
added again as a direct 1PI \(p_j\) contact.

\subsection*{Coefficient lock}

For a canonically normalized image oscillator, the most general single-pole
quadratic form with the same tensor structure is
\begin{equation}
\begin{split}
  S_j^{(2)}
  ={}&\langle p_j,Bp_j\rangle
  +\langle Q_j,(B+\alpha\mu^2A_jA_j^\dagger)Q_j\rangle\\
  &-\gamma\mu^2\langle Q_j,A_jp_j\rangle
  -\gamma\mu^2\langle A_jp_j,Q_j\rangle .
  \label{eq:general-osc}
\end{split}
\end{equation}
Eliminating \(Q_j\) gives
\begin{equation}
  \Gamma_j(B)
  =
  B-\frac{\gamma^2\mu^4J_j}{B+\alpha\mu^2J_j}.
  \label{eq:general-gamma}
\end{equation}
With \(\mu_{\eff}^2=\alpha\mu^2\) and \(X=B/\mu_{\eff}^2\), the pole equation is
\begin{equation}
  X=\frac{r^2J_j}{X+J_j},
  \qquad r=\frac{\gamma}{\alpha}.
  \label{eq:r-ratio}
\end{equation}
The de Vries pole requires \(r=1\).

The local square
\begin{equation}
  \alpha\mu^2\|A_j^\dagger Q_j-p_j\|^2
  -\alpha\mu^2\|p_j\|^2
  \label{eq:single-square}
\end{equation}
expands to \eqref{eq:general-osc} with \(\gamma=\alpha\).  Therefore a single
current-displacement square gives the needed coefficient lock; a generic
oscillator does not.

DBI/Stueckelberg displacement terms have this one-square form, so they provide
the right coefficient mechanism.  The present D8 normal displacement is only a
scalar on \(\WW_9\), however; it is not a field in
\(\Im A_j\subset V_j\otimes\RR^3\).  Thus the finite branch still requires an
\(\Im A_j\) vector-bundle displacement supplied by the D8/open-current
interface.

The mixed Chan--Paton boundary-changing sector supplies a sharper candidate.
For \(j>0\),
\begin{equation}
  H^0(\CP^1,\operatorname{Hom}(\OO(0),\OO(2j)))=V_j,
  \qquad
  U_j=\frac{1}{\sqrt{J_j}}A_j:V_j\to\Im A_j
  \label{eq:hom-u}
\end{equation}
is an isometry.  Thus a mixed zero mode \(y_j\in V_j\) gives a canonical image
field
\begin{equation}
  Q_j=U_jy_j=\frac{1}{\sqrt{J_j}}A_jy_j,
  \qquad
  A_j^\dagger Q_j=\sqrt{J_j}\,y_j .
  \label{eq:hom-q}
\end{equation}
The desired current-displacement square is equivalent to
\begin{equation}
  \alpha\mu^2\|\sqrt{J_j}\,y_j-p_j\|^2
  -\alpha\mu^2\|p_j\|^2 .
  \label{eq:hom-square}
\end{equation}
This passes the field-content test but not the local-action test.  A generic
Hom-mode quadratic form
\begin{equation}
\begin{split}
  S_y^{(2)}
  ={}&\langle p_j,Bp_j\rangle
  +Z_y\langle y_j,By_j\rangle
  +c_m\mu^2J_j\langle y_j,y_j\rangle\\
  &-c_g\mu^2\sqrt{J_j}
  \bigl(\langle y_j,p_j\rangle+\langle p_j,y_j\rangle\bigr)
\end{split}
\end{equation}
gives
\begin{equation}
  \Gamma_j(B)=B-\frac{c_g^2\mu^4J_j}{Z_yB+c_m\mu^2J_j},
  \qquad
  r=\frac{c_g\sqrt{Z_y}}{c_m}.
  \label{eq:hom-r}
\end{equation}
The de Vries pole requires \(c_m=c_g\sqrt{Z_y}\).  The mixed zero-mode
statement supplies the bundle; it does not by itself supply this coefficient
lock.

The current source package contains no term of the form \eqref{eq:hom-square}.  The Payen
endpoint action gives the representation, Wilson loop, and charge insertions
\(T_a,T_aT_b\).  It is first order on the boundary and contains no propagating
Hom-mode potential.  The D8 DBI/WZ package supplies geometric brane norms and
RR charge data, but no term identifying the mixed Hom mode with the mismatch
\(\sqrt{J_j}y_j-p_j\).  With the current action, the Hom sector therefore has
the generic ratio \eqref{eq:hom-r}.

\section{Finite branch and tower branch}

The residue calculation separates two branches.

\begin{center}
\small
\begin{tabular}{@{}>{\raggedright\arraybackslash}p{0.24\linewidth}
>{\raggedright\arraybackslash}p{0.32\linewidth}
>{\raggedright\arraybackslash}p{0.32\linewidth}@{}}
\hline
Branch & Spectrum rule & Phenomenological signal\\
\hline
Finite current branch &
The level-two endpoint/current sector keeps \(j=0,\frac12,1\).  The Feshbach
pole uses \eqref{eq:inverse-image}. &
No \(j=3/2\) state in the retained weak boundary spectrum.\\
Tower/moduli branch &
The rotor or neighboring moduli slice keeps higher \(j\).  The natural gap is
proportional to \(j(j+1)\). &
The first extra positive branch is \(j=3/2\), near the low-mass diphoton window.\\
\hline
\end{tabular}
\end{center}

For \(j=3/2\), \(J=15/4\) and
\begin{equation}
  X_+\!\left(\frac{15}{4}\right)=0.8204823316059779.
  \label{eq:x32}
\end{equation}
Fixing the common scale by \(M_Z=91.1880\,{\rm GeV}\) gives
\begin{equation}
  M_{3/2}
  =
  M_Z\sqrt{\frac{X_+(15/4)}{X_+(2)}}
  =
  96.53875533197167\,{\rm GeV}.
  \label{eq:m32}
\end{equation}
Thus a tower/moduli branch places a neutral \(j=3/2\) mass inside the CMS Run 2
low-mass diphoton search range and \(1.14\,{\rm GeV}\) above the reported
\(95.4\,{\rm GeV}\) local feature.

\section{Extra-mode gap check}

The finite image operator is safe only on the image itself.  If the field is an
unprojected element of \(V_j\otimes\mathbb R^3\), then
\[
  V_j\otimes\mathbb R^3=\Im A_j\oplus(\Im A_j)^\perp .
\]
On \(\Im A_j\), \(A_jA_j^\dagger=J_j\).  On the complement,
\[
  A_j^\dagger q_\perp=0,
  \qquad
  A_jA_j^\dagger q_\perp=0.
\]
Thus the quadratic operator \(B+\mu^2A_jA_j^\dagger\) gives denominator \(B\)
on \(q_\perp\).  The complement is massless unless the action supplies the
projection \(q_j=P_jq_j\) or a positive complement mass.

The mixed Hom field avoids this complement because
\[
  y_j\in H^0(\CP^1,\operatorname{Hom}(\OO(0),\OO(2j)))=V_j,
  \qquad
  Q_j=\frac{1}{\sqrt{J_j}}A_jy_j\in\Im A_j.
\]
This fixes the finite field content, but the current D8/Payen/RR package
contains no one-square Hom action.

The rotor branch has the opposite problem.  With the \(Z\) fixing the scale,
the next tower masses are
\[
\begin{array}{c|c}
j & M_j\ {\rm GeV}\\
\hline
3/2 & 96.5388\\
2 & 99.5795\\
5/2 & 101.4539\\
3 & 102.6807
\end{array}
\]
and the large-\(j\) limit is \(106.57\,{\rm GeV}\).  The rotor branch is
therefore not a high-gap branch.

\section{Coupled selector potential}

The boundary scale can be coupled to closed-string moduli by writing
\(\mu_B=\mu_B(\rho,\tau)\).  The Type IIA scaling potential is
\begin{equation}
\begin{split}
  V_{\rm bulk}
  ={}&-B_R\tau^{-2}\rho^{-2}
  +B_H\tau^{-2}\rho^{-6}
  +B_0\tau^{-4}\rho^6
  +B_2\tau^{-4}\rho^2\\
  &+B_4\tau^{-4}\rho^{-2}
  +B_6\tau^{-4}\rho^{-6}
  +B_{\rm loc}\tau^{-3}.
  \label{eq:vbulk}
\end{split}
\end{equation}
Set
\begin{equation}
  J_v=1+\sqrt3,
  \qquad
  J_H=\frac{3+\sqrt{57}}{8},
  \qquad
  \lambda_{\rm dV}
  =
  \frac{J_H}{4J_v}
  =
  0.12067210653215474.
  \label{eq:scalar-targets}
\end{equation}
The coupled selector potential is
\begin{equation}
\begin{split}
  V_{\rm sel}
  ={}&C_{\rm cur}(\rho,\tau)(\kappa_{\rm cur}-1)^2\\
  &+C_v(\rho,\tau)
  \left(v_Q^2-\mu_B(\rho,\tau)^2J_v\right)^2\\
  &+C_H(\rho,\tau)(\lambda_H-\lambda_{\rm dV})^2,
  \qquad C_{\rm cur},C_v,C_H>0.
  \label{eq:vsel}
\end{split}
\end{equation}
At the radial minimum \(h^2=2v_Q^2\), stationarity in the selector variables
gives
\begin{equation}
  \kappa_{\rm cur}=1,
  \qquad
  v_Q^2=\mu_B(\rho,\tau)^2(1+\sqrt3),
  \qquad
  \lambda_H=\lambda_{\rm dV}.
  \label{eq:selector}
\end{equation}
With \(v_{\rm EW}=\sqrt2\,v_Q\),
\begin{equation}
  M_H^2=2\lambda_Hv_{\rm EW}^2.
  \label{eq:higgs}
\end{equation}
On the selector branch the squared selector terms vanish, so the \(\rho,\tau\)
stationarity equations reduce to the bulk coefficient balances from
\eqref{eq:vbulk}.

The selector terms do not select \(\rho,\tau\) on that branch.  If
\[
  S_{\rm cur}=\kappa_{\rm cur}-1,\quad
  S_v=v_Q^2-\mu_B(\rho,\tau)^2J_v,\quad
  S_H=\lambda_H-\lambda_{\rm dV},
\]
then \(V_{\rm sel}=\sum_i C_i(\rho,\tau)S_i^2\), and at \(S_i=0\)
\[
  \partial_\rho V_{\rm sel}
  =
  \sum_i(\partial_\rho C_i)S_i^2
  +2\sum_iC_iS_i\partial_\rho S_i
  =0
\]
with the same statement for \(\partial_\tau\).  The bulk equations are
therefore the only \(\rho,\tau\) stationarity equations in this written
potential.

Those bulk equations are two linear constraints on the seven coefficients
\[
  B_R,B_H,B_0,B_2,B_4,B_6,B_{\rm loc}.
\]
For instance, in the reduced coefficient slice
\[
  B_H=B_0=B_2=B_6=0,
\]
stationarity at any chosen \(\rho\) is obtained by
\[
  B_4=B_R\tau^2,
  \qquad
  B_{\rm loc}=-\frac{2}{3}B_R\tau\rho^{-2}.
\]
At \(\rho_*=9.463469757158872\), this gives
\[
  B_{\rm loc}=-0.007444027070759169\,B_R\tau .
\]
Thus the written vacuum potential supplies a coefficient-balance
parametrization, not a selection of \(\rho_*\).  The source-balance obstruction
below also removes the current D8/O8 localized source package as a global
completion for \(B_{\rm loc}\).

\section{Neutral scalar Hessian}

The scalar constants in \eqref{eq:scalar-targets} are the absolute values of
the negative root of the same quadratic as \eqref{eq:xplus}:
\begin{equation}
  |X_-(2)|=1+\sqrt3=J_v,
  \qquad
  |X_-(3/4)|=\frac{3+\sqrt{57}}{8}=J_H .
  \label{eq:xminus-scalar}
\end{equation}
This extends the de Vries algebra to the scalar slice.  The local scalar
Hessian test asks whether the boundary action forces this extension.

For a canonically normalized neutral radial field \(h\), write
\begin{equation}
  V_h(h)=a(\rho,\tau)h^2+b(\rho,\tau)h^4+V_0(\rho,\tau),
  \qquad b>0 .
  \label{eq:vh-general}
\end{equation}
At a nonzero stationary point,
\begin{equation}
  h_0^2=-\frac{a}{2b},
  \qquad
  M_h^2=V_h''(h_0)=8bh_0^2=-4a .
  \label{eq:hessian-general}
\end{equation}
The selector normalization is \(h_0^2=v_{\rm EW}^2=2v_Q^2\).  If
\(v_Q^2=J_v\mu_B^2\), then
\begin{equation}
  \frac{M_h^2}{\mu_B^2}=16bJ_v .
  \label{eq:mh-free-b}
\end{equation}
Equivalently,
\begin{equation}
  V_h=\frac{\lambda_H}{4}(h^2-2v_Q^2)^2
  \quad\Rightarrow\quad
  M_h^2=4\lambda_Hv_Q^2 .
  \label{eq:vh-lambda}
\end{equation}
The scalar target \(M_h^2=J_H\mu_B^2\) is obtained by setting
\begin{equation}
  \lambda_H=\frac{J_H}{4J_v}.
  \label{eq:lambda-target-repeat}
\end{equation}
Thus the scalar formula is internally consistent, but the coefficient
\(b\), or equivalently \(\lambda_H\), is a free scalar coefficient in the
present local potential.

Adding neutral scalar mixing does not change this rank count.  For neutral
variables \(\Phi^A=(h,\phi^i)\), the physical scalar masses are eigenvalues of
\begin{equation}
  H_{AB}=\frac{\partial^2V}{\partial\Phi^A\partial\Phi^B}.
  \label{eq:neutral-hessian}
\end{equation}
Requiring one eigenvalue to be \(J_H\mu_B^2\) is one spectral condition on the
entries of \(H_{AB}\).  Positivity restricts signs, but it does not select the
square-root value \(J_H\).  The current scalar sector therefore realizes the
negative branch as a compatible target, not as a forced Hessian eigenvalue.

\section{Source balance}

The D8/O8 WZ bookkeeping gives a compact source table.  Let
\begin{equation}
  \frac{F_Q}{2\pi}=f_mh_Q,\qquad f_m=\frac{2m+13}{2}.
\end{equation}
The D8 stack contributes
\begin{equation}
  Q_{D8}=N_8,\qquad
  Q_{D6}=f_mN_8h_Q,\qquad
  Q_{D4}=\frac{2(2m+13)^2-1}{16}N_8h_Q^2.
\end{equation}
The local O8 contribution is
\begin{equation}
  Q_{O8}=-16-\frac12h_Q^2.
\end{equation}
Spin-c integrality requires \(N_8\in16\mathbb Z\).  In a two-endpoint type I'
bookkeeping with total stack \(N_{8,L}+N_{8,R}=32\), the total eight-charge
cancels, but the induced lower vector is
\begin{equation}
  Q_{\rm lower}^{\rm total}
  =
  32f_mh_Q+
  \left[2(2(2m+13)^2-1)-1\right]h_Q^2.
  \label{eq:qlower}
\end{equation}
For the minimal spin-c lift,
\begin{equation}
  m=0:\qquad Q_{\rm lower}^{\rm total}=208h_Q+673h_Q^2.
  \label{eq:qlower0}
\end{equation}
Thus the current source package does not pass the global tadpole test.  The
eight-charge can be balanced, but no current bulk flux, image sector, K-theory
identification, or extra localized source cancels \eqref{eq:qlower0}.

More generally, let an added lower-source sector contribute
\begin{equation}
  Q_{\rm new}=\alpha h_Q+\beta h_Q^2 .
  \label{eq:qnew}
\end{equation}
The lower Bianchi equations require
\begin{equation}
  \alpha=-\frac{N_LM_L+N_RM_R}{2},
  \label{eq:alpha-target}
\end{equation}
and
\begin{equation}
  \beta=
  -\left[
  \frac{(2M_L^2-1)N_L+(2M_R^2-1)N_R}{16}-1
  \right].
  \label{eq:beta-target}
\end{equation}
For an unequal endpoint stack this is
\begin{equation}
  Q_{\rm new}=-16M h_Q-(4M^2-3)h_Q^2 ,
  \qquad M\in2\mathbb Z+1.
  \label{eq:unequal-source-target}
\end{equation}
For a balanced endpoint stack with \(M_R=-M_L\), the D6 entry vanishes and
\begin{equation}
  Q_{\rm new}=-(4M_L^2-3)h_Q^2 .
  \label{eq:balanced-source-target}
\end{equation}
The minimal balanced choice \(M_L=\pm1\) therefore needs
\begin{equation}
  Q_{\rm new}=-h_Q^2 .
  \label{eq:min-source-target}
\end{equation}

This obstruction is not removed by changing the two endpoint spin-c lifts
inside the present type I' data.  Write \(M_i=2m_i+13\).  Unequal stack
distributions leave
\begin{equation}
  Q_6=16M_i h_Q,
\end{equation}
which cannot vanish because \(M_i\) is odd.  Balanced stacks can cancel D6
only when \(M_R=-M_L\).  Then
\begin{equation}
  Q_4=(4M_L^2-3)h_Q^2,
\end{equation}
whose minimum is \(h_Q^2\).  Thus the endpoint data alone cannot cancel all
lower charges.

Closed \(H\)-flux does not remove the obstruction on the selected spaces.  The
product and flag cohomology rings have only even generators, hence
\begin{equation}
  H^3=H^5=0 .
\end{equation}
With no NS5 source, \([H]=0\), so the cohomology classes \([HF_0]\) and
\([HF_2]\) vanish.  The existing NSNS flux entry in the potential scaffold is
therefore not a cancellation of the lower source vector.

Nor can the current lower vector be removed by a bare K-theory identification.
For the product and flag spaces the cohomology is torsion-free and
concentrated in even degree, so the rational Chern character detects the
lower charge entries.  A nonzero rational \(h_Q\) or \(h_Q^2\) component is not
K-theory equivalent to zero in the charge group used by the WZ action.

The lower-only source requirement can be isolated.  In the same-lift source
table, an added image sector with no D8 charge would need
\begin{equation}
  Q_{\rm img}=-208h_Q-673h_Q^2.
\end{equation}
The balanced opposite-lift endpoint package cancels the D6 part first and
leaves the residual D4 source
\begin{equation}
  Q_4=(4M_L^2-3)h_Q^2,
  \qquad M_L\in2\mathbb Z+1.
\end{equation}
The smallest residual is \(h_Q^2\).  Hence the minimal lower-only completion is
a pure negative D4 image/source
\begin{equation}
  Q_{\rm img}=-h_Q^2.
\end{equation}
A literal anti-D4 or an O4-like image source has this charge type, but either
choice is an enlargement of the source package and brings additional localized
tension, backreaction, and potential coefficients.

The O4 normalization is favorable.  The orientifold WZ coefficient
\(-2^{p-5}\) gives an \(O4^-\) charge \(-1\) in the doubled type I' convention.
Thus an \(O4^-\) fixed component on \(\{pt\}_{\CP^2}\times S^1_Q\) would supply
\(-h_Q^2\).  The product fixed-locus geometry does not give this component by
itself.  A holomorphic projective involution on \(\CP^2\) fixes
\(\CP^0\sqcup\CP^1\), so with the \(RP^1\) fixed circle in \(\CP^1\) it
produces O4 and O6 supports.  Standard complex conjugation fixes \(RP^2\) and
therefore gives O6-type support after the \(\CP^1\) fixed circle is included.
The O4 charge match therefore requires an enlarged O4/O6 source table rather
than closing the present D8/O8 package.

The enlarged holomorphic package still does not close the source table.  In the
same doubled convention the companion \(O6^-\) carries \(-4h_Q\), so one
product involution contributes
\begin{equation}
  Q_{O4/O6}=-4h_Q-h_Q^2 .
\end{equation}
The D8 endpoint D6 charge is
\begin{equation}
  Q_6^{D8}=\frac{N_LM_L+N_RM_R}{2}h_Q,
  \qquad N_i\in\{0,16,32\},\quad M_i\in2\mathbb Z+1,
\end{equation}
hence \(Q_6^{D8}\in16\mathbb Z\,h_Q\).  It cannot cancel the single O6 unit.
With \(k\) identical O4/O6 packages, D6 cancellation forces \(k\) to be a
multiple of four, while the balanced endpoint D4 coefficient \(4M^2-3\) is
\(1\) modulo \(8\).  Thus the O4/O6 enlargement does not close Milestone 6
inside the current endpoint lattice.

The source branch is therefore closed for the present compactification
ansatz.  A standalone anti-D4 would cancel the minimal cohomology class, but it
adds localized tension, an open sector, and backreaction not included in the
boundary-current action.  An O4-only repair would require a different
orientifold or topology.  A D6-bearing repair adds a new source sector and
changes the Bianchi table.  Each option is a new compactification ansatz, not a
completion of the current D8/O8 route.

\begin{proposition}[Source-balance obstruction]
For the current D8/O8 endpoint package on the selected product or flag
topology, with closed \(H\), no NS5 source, rational WZ/K-theory charge
detection, and only the tested product O4/O6 enlargement, the lower Bianchi
classes cannot vanish.
\end{proposition}

\begin{proof}
Endpoint D8/O8 data alone fail as follows.  Unequal endpoint stacks leave
\(Q_6=16M_i h_Q\), which is nonzero for odd \(M_i\).  Balanced stacks cancel
D6 only with \(M_R=-M_L\), and then leave
\((4M_L^2-3)h_Q^2\), which is nonzero.  Closed \(H\) has no cohomology class on
the selected spaces because \(H^3=H^5=0\), so the \(HF_0\) and \(HF_2\) terms
cannot cancel these classes.  The rational Chern character detects the
nonzero \(h_Q\) and \(h_Q^2\) entries on the torsion-free even cohomology
spaces, so bare K-theory cannot identify the lower vector with zero.  The
tested holomorphic O4/O6 package contributes \(-4k h_Q-kh_Q^2\).  D6
cancellation forces \(k\) to be a multiple of four, but the balanced endpoint
D4 coefficient is \(1\) modulo \(8\), and the unequal endpoint D4 equation
\(4M_i^2-3=4M_i\) has no integer solution.  Thus all tested cancellation
routes fail under the stated assumptions.
\end{proof}

\begin{proposition}
Under the scoped assumptions above, the de Vries weak-angle relation is not
forced by the D10 dynamics.  The construction is a
compatibility framework: it contains the D10 carrier/probe geometry, the
endpoint Casimirs, the finite image-channel algebra, and the full-current
metric branch, but at least one mandatory pass condition fails on every route
that would turn the algebraic match into a prediction.
\end{proposition}

\begin{proof}
Milestone 1 fails as a forcing statement because the scoped action does not
supply the one-square Hom/image term needed to lock the Feshbach residue.
Milestone 2 selects the finite set only on an imposed finite branch; Borel-Weil
alone gives a tower.  Milestone 3 gives \(\kappa_{\rm cur}=1\) only on the
full-current branch, while the scoped source package does not select that
branch.  Milestones 4 and 5 leave the bulk coefficients, boundary scale, and
scalar Hessian as coefficient targets.  Milestone 6 fails global source balance
by the preceding proposition and closes the current source branch.  Milestone 7
therefore has no derived source/vacuum system from which to compute thresholds,
scheme conversion, or branching.  Thus the de Vries number is compatible with
the scoped route but is not a prediction of it.
\end{proof}

\section{Status of the construction}

\begin{proposition}[Current result]
The D10 boundary-current construction gives a coherent compatibility framework
for the de Vries weak-angle number.  The canonical current metric gives
\(\kappa_{\rm cur}=1\) on the full-current branch, but the written D8/Payen/RR
source package contains no term fixing that branch.  The finite image-channel
Feshbach algebra gives the de Vries pole only when the open-channel contact has
unit residue.  The raw Wilson contact has residue \(J_j\), so ordinary Wilson
expansion plus image projection does not force the finite pole when treated as
a separate direct 1PI term.  In the retained-field action the Wilson second
derivative is a connected response represented by \(q_j\) exchange.  The local
oscillator must also satisfy the coefficient lock \(r=1\) in
\eqref{eq:r-ratio}.  The mixed Hom sector supplies the right finite image
field through \eqref{eq:hom-q}, but Payen plus the D8 DBI/WZ source package
contains no one-square Hom action \eqref{eq:hom-square}.  The source
table leaves the lower charge vector \eqref{eq:qlower0} uncancelled; after the
best endpoint-lift cancellation it still needs a negative D4 image/source
\(-h_Q^2\).  The \(O4^-\) charge magnitude matches this unit, but the product
involutions add O6-type support.  The resulting O6 charge is incompatible with
the current D8 endpoint lattice, so the source package is not closed.  The
scoped source-balance obstruction above rules out the current D8/O8 package as
a global compactification proof.  The extra-mode check also fails for the
current action: an unprojected image field has a massless complement, while
the rotor branch keeps a light tower.  The vacuum rank test shows that the
written selector and bulk potential do not select \(\rho_*\).  The scalar
Hessian rank test shows that the negative-branch scalar constants are fitted by
free scalar coefficients in the present action.  The replacement-source
equations \eqref{eq:alpha-target}--\eqref{eq:min-source-target} reduce the
minimal source repair to \(-h_Q^2\), which is not supplied by the current
package.  Accepting that repair changes the source content or topology, so the
current compactification branch is not a proof branch.
\end{proposition}

\begin{proof}
The compatibility framework is the collection of \eqref{eq:d10},
\eqref{eq:w9}, \eqref{eq:vj}, \eqref{eq:qeta}, \eqref{eq:squad},
\eqref{eq:kappaone}, and \eqref{eq:vsel}.  The Feshbach pole follows from
\eqref{eq:gamma}.  The current metric follows from \eqref{eq:metric} and
\eqref{eq:frame}.  The residue mismatch follows from comparing
\eqref{eq:creq} with \eqref{eq:cwilson}.
\end{proof}

\begin{thebibliography}{9}

\bibitem{Witten1995}
E. Witten,
``String theory dynamics in various dimensions,''
Nucl. Phys. B \textbf{443}, 85--126 (1995).

\bibitem{CMSHIG20002}
CMS Collaboration,
``Search for a standard model-like Higgs boson in the mass range between 70 and
110 GeV in the diphoton final state in proton-proton collisions at
\(\sqrt{s}=13\,{\rm TeV}\),''
Phys. Lett. B \textbf{793}, 320--347 (2019);
CMS-HIG-20-002 public result page.

\bibitem{PDG2025}
Particle Data Group,
``Gauge and Higgs boson summary tables,''
Review of Particle Physics (2025).

\end{thebibliography}

\end{document}
